% Short report on the black-hole and neutron-star production runs of 2026-09-17/18.
%
% Build:  cd Report && latexmk -pdf populations_report.tex
% Figures are pulled from ../figures/prod_20260918/ via \graphicspath below; they are the
% exact files the analysis scripts wrote, not copies, so this document cannot drift from them.
\documentclass[english,10pt,a4paper]{article}
\usepackage[T1]{fontenc}
\usepackage[margin=2.4cm]{geometry}
\usepackage{graphicx}
\graphicspath{{./}{../figures/prod_20260918/}}
\usepackage{mathtools}
\usepackage{amssymb}
\usepackage{booktabs}
\usepackage{xcolor}
\usepackage[colorlinks=true,allcolors=blue]{hyperref}
\usepackage{caption}

\newcommand{\tE}{t_{\rm E}}
\newcommand{\piE}{\pi_{\rm E}}
\newcommand{\tetE}{\theta_{\rm E}}
\newcommand{\Ml}{M_{\rm L}}
\newcommand{\Msun}{M_\odot}
\newcommand{\Neff}{N_{\rm eff}}

\title{Compact objects in the Roman--Rubin joint microlensing forecast:\\
       black holes and neutron stars}
\author{Ali Salehi}
\date{\today}

\begin{document}
\maketitle

\begin{abstract}
We report two production runs of the joint Roman\,+\,Rubin Galactic-bulge microlensing
simulation, one for a population of isolated black holes with a log-uniform mass function over
$3$--$1000\,\Msun$, and one for neutron stars drawn from a Gaussian at $1.35\pm0.15\,\Msun$
truncated to $1.10$--$2.20\,\Msun$. Each run scans 1829 sightlines over $\sim\!68\,\mathrm{deg}^2$
and characterises every detected event three times -- with Rubin alone, with Roman alone, and
with the two combined -- through independent Fisher matrices. All pooled quantities are
event-rate weighted and quoted with the Kish effective sample size. We find that the joint fit
improves the timescale by a factor $6$--$9$ over Rubin alone and $\sim\!1.35$ over Roman
alone, and that it barely improves $\tetE$ at all, which is set by astrometry and therefore by
Roman. The centroid shift exceeds Roman's per-exposure astrometric floor for $90\%$ of
black-hole events and $0.6\%$ of neutron-star events. Satellite parallax from L2 is found to be
almost worthless for both populations. Following Sajadian \& Makler (2026), we compute the
probability that the two lensing-induced images are separately resolvable, and find it to be
dominated by the assumed resolution criterion rather than by the lens population.
\end{abstract}

\section{Simulation and runs}

The simulator draws source stars and lenses along each sightline from Besan\c{c}on/CMD
population models with three-dimensional dust extinction, generates per-band light curves
against the real Rubin visit list and the published Roman GBTDS cadence, applies a Monte Carlo
detection test, and forecasts parameter precision with a Fisher matrix. The photometric matrix
has nine parameters $(u_0,\tE,f_{b,0},\piE,\xi,t_0,m_{b,0},f_{b,1},m_{b,1})$ and the astrometric
matrix four $(\tetE,\mu_{s,1},\mu_{s,2},\piE)$; which parameters are free is decided per event
from the epochs each survey actually contributed. Crucially there is one such pair of matrices
\emph{per survey partition} -- Rubin-only, Roman-only, and joint -- so that ``what the
combination adds'' is measured rather than assumed.

Both runs used identical configuration apart from the lens mass function
(\texttt{--events 300 --lenses 50 --stride-roman 5 --pair-satellite}), a fixed detection
threshold $\Delta\chi^2 \ge 500$, and the same RNG seed. The stratified grid places sightlines
every $0.1^\circ$ inside Roman's footprint and every $0.2^\circ$ outside it, which oversamples
the footprint by design; the per-sightline sky area $w_{\rm area}$ undoes that in every pooled
statistic. Table~\ref{tab:runs} summarises the outcome.

\begin{table}[h]
\centering
\caption{Run configuration and outcome. ``Barren'' sightlines drew stars but characterised
nothing; they run to the full draw cap and so dominate the raw row count while contributing no
detections whatsoever. The weightable sample is the one every number in this report is computed
from.}
\label{tab:runs}
\begin{tabular}{lrr}
\toprule
 & black holes & neutron stars \\
\midrule
mass function            & log-uniform          & Gaussian (\"Ozel \& Freire 2016) \\
mass range [$\Msun$]     & $3$--$1000$          & $1.10$--$2.20$ \\
sightlines scanned       & 1829                 & 1829 \\
\quad aggregated         & 1612                 & 1612 \\
\quad no coverage        & 140                  & 140 \\
\quad barren             & 77                   & 77 \\
\quad hit draw cap       & 77                   & 77 \\
rows written             & 4\,679\,203          & 5\,012\,349 \\
weightable draws         & 829\,202             & 1\,162\,348 \\
detected events          & 110\,144             & 93\,685 \\
$\Neff$ (weighted)       & 143\,499             & 296\,562 \\
CPU time                 & $\sim\!13.5$\,h      & $\sim\!13.5$\,h \\
\bottomrule
\end{tabular}
\end{table}

\paragraph{Weighting.} The Monte Carlo draws the lens \emph{distance} in proportion to the event
rate but the lens \emph{mass} from the number density and the velocities from plain Gaussians,
so it is missing the rate's $\sqrt{\Ml}\,v_t$ factor. Every pooled fraction, median or ratio in
this report therefore carries the per-event weight
$W \propto w_{\rm area}\,N_\star/n_{\rm sim}\,\sqrt{\Ml}\,v_t\,Z(D_s)$, and every weighted number
is quoted with $\Neff = (\sum w)^2/\sum w^2$, because the weight costs precision. Validation:
weighting all draws of an earlier bulge run reproduces the OGLE-IV efficiency-corrected mean
timescale of $22$\,d (we obtain $24.0$\,d; unweighted the same sample gives $56.2$\,d).

\section{What each survey contributes}

\begin{figure}[h]
\centering
\includegraphics[width=\textwidth]{p6_synergy.png }
\caption{(a) Weighted share of detected events found by each survey. (b) Median ratio of the
joint to the single-survey timescale error, \emph{restricted to events that both surveys
characterised independently}. That restriction is essential: Roman contributes epochs to only a
few per cent of draws, so pooled over all characterised events the ``joint'' fit is literally the
Rubin fit for the great majority and the ratio is $1.000$ by construction.}
\label{fig:synergy}
\end{figure}

Rubin finds the overwhelming majority of events in both populations -- it covers the whole
$68\,\mathrm{deg}^2$ scan, while Roman's footprint is $\sim\!2\%$ of it -- but the events the two
share are the ones that are well measured. Table~\ref{tab:synergy} gives the paired comparison.

\begin{table}[h]
\centering
\caption{Detection provenance and the joint-fit gain. The $\sigma$ ratios are weighted medians
over events characterised independently by \emph{both} surveys (15\,849 black-hole and 12\,569
neutron-star events); values below unity mean the joint fit is better.}
\label{tab:synergy}
\begin{tabular}{lrrrr}
\toprule
 & \multicolumn{2}{c}{black holes} & \multicolumn{2}{c}{neutron stars} \\
\cmidrule(lr){2-3}\cmidrule(lr){4-5}
 & vs.\ Rubin & vs.\ Roman & vs.\ Rubin & vs.\ Roman \\
\midrule
detected by Rubin only [\%] & \multicolumn{2}{c}{89.4} & \multicolumn{2}{c}{90.0} \\
detected by Roman only [\%] & \multicolumn{2}{c}{5.3}  & \multicolumn{2}{c}{6.3} \\
detected by both [\%]       & \multicolumn{2}{c}{5.2}  & \multicolumn{2}{c}{3.6} \\
\midrule
median $\sigma(\tE)$ ratio    & 0.165 & 0.733 & 0.116 & 0.742 \\
median $\sigma(\piE)$ ratio   & 0.108 & 0.826 & 0.114 & 0.798 \\
median $\sigma(\tetE)$ ratio  & 0.187 & 0.979 & 0.152 & 0.975 \\
\bottomrule
\end{tabular}
\end{table}

The asymmetry is the result. Roman's dense cadence sharpens what Rubin alone could do on the
timescale by a factor $6$--$9$ and on the parallax by $\sim\!9$. Rubin's decade-long baseline
sharpens Roman by a more modest $\sim\!1.25$--$1.35$. But for $\tetE$ the ratio against Roman
alone is $0.975$--$0.979$: Rubin adds essentially nothing, because $\tetE$ is an astrometric
quantity and Roman is the only astrometric instrument in this pairing. Any science case resting
on $\tetE$ -- which includes the lens mass, through $\Ml = \tetE/(\kappa\piE)$ -- is a Roman
case, not a joint one.

\section{Fisher forecast precision}

\begin{figure}[h]
\centering
\includegraphics[width=\textwidth]{p7_precision.png}
\caption{Weighted cumulative distribution of the fractional parameter error, per survey and
population. \textbf{Each curve is over its own characterised set}, not a common one: Roman
characterises only inside its footprint, where coverage is good, while the joint curve also
contains every Rubin-only event. That is why Roman's $\tetE$ curve lies above the joint curve in
panel (c), which is a difference of samples and not a violation of
$\sigma_{\rm joint}\le\sigma_{\rm single}$.}
\label{fig:precision}
\end{figure}

\begin{table}[h]
\centering
\caption{Weighted fraction of characterised events measured to better than $10\%$. Read down a
column, not across: each entry has its own denominator (see Fig.~\ref{fig:precision}).}
\label{tab:tenpercent}
\begin{tabular}{lrrrrrr}
\toprule
 & \multicolumn{3}{c}{black holes [\%]} & \multicolumn{3}{c}{neutron stars [\%]} \\
\cmidrule(lr){2-4}\cmidrule(lr){5-7}
 & joint & Rubin & Roman & joint & Rubin & Roman \\
\midrule
$\tE$   & 19.77 & 17.66 & 18.61 & 7.71 & 6.12 & 10.51 \\
$\piE$  &  3.13 &  2.14 &  7.71 & 1.50 & 1.08 &  2.91 \\
$\tetE$ & 43.87 & 41.11 & 81.41 & 6.03 & 0.37 & 42.28 \\
$\Ml$   &  1.50 &  0.44 &  7.07 & 0.34 & 0.01 &  1.85 \\
\bottomrule
\end{tabular}
\end{table}

The lens mass is the hardest quantity in the table, as expected: it needs both $\tetE$ and
$\piE$, and $\piE$ is the weaker of the two. Within Roman's footprint, where the paired
comparison is meaningful, the median $\sigma(\Ml)/\Ml$ improves by a factor $0.121$ (black
holes) and $0.138$ (neutron stars) relative to Rubin alone, but only to $0.862$ and $0.806$
relative to Roman alone -- the same asymmetry as $\tetE$, and for the same reason. Note the
footprint sample is much smaller than the full scan: $\Neff = 5471$ and $5372$ respectively, so
these are the least precise numbers in this report.

\begin{figure}[h]
\centering
\includegraphics[width=\textwidth]{p7_precision_te.png}
\caption{Median fractional error against the event timescale. The two surveys fail at opposite
ends: Rubin's $\sim\!3$-day cadence under-resolves short events, while Roman's $72$-day seasons
cannot contain long ones, which is where Rubin's decade of coverage matters.}
\label{fig:precision_te}
\end{figure}

\section{Astrometric microlensing}

\begin{figure}[h]
\centering
\includegraphics[width=\textwidth]{p6_astrometry.png}
\caption{(a) Weighted distribution of the maximum centroid shift reached by each event. (b) The
blend-diluted shift against lens mass, with the weighted interquartile band; the dotted line is
Roman's $1.1$\,mas per-exposure astrometric floor.}
\label{fig:astrometry}
\end{figure}

The centroid of the unresolved image pair is displaced by
$\delta\theta(u) = \tetE\,u/(u^2+2)$, which peaks at $u=\sqrt{2}$ and \emph{not} at closest
approach; an event whose trajectory crosses $u=\sqrt{2}$ therefore reaches $\tetE/\sqrt{8}$.
Because $\tetE \propto \sqrt{\Ml}$, this is the channel where the two populations diverge most
sharply (Table~\ref{tab:ast}): black holes cross Roman's per-exposure floor at
$\Ml \approx 20\,\Msun$ and $90\%$ of their detections exceed it, while neutron stars sit an
order of magnitude below it. Blending is a small correction here -- the Roman-band source flux
fraction averages $0.95$ -- but it is included rather than assumed negligible.

\begin{table}[h]
\centering
\caption{Astrometric shift, satellite parallax, and image resolvability. The resolution
probability is the weighted fraction of that survey's detections having at least three epochs at
which both images are separately detectable and separated by more than the stated bar.}
\label{tab:ast}
\begin{tabular}{lrr}
\toprule
 & black holes & neutron stars \\
\midrule
median max.\ centroid shift [mas]      & 3.61  & 0.265 \\
95th percentile [mas]                  & 9.75  & 0.610 \\
fraction above $1.1$\,mas floor [\%]   & 90.2  & 0.57 \\
\midrule
satellite parallax, median gain        & 1.000 & 1.000 \\
\quad fraction gaining $>1.1\times$ [\%] & 0.003 & 0.200 \\
\quad fraction gaining $>2\times$ [\%]   & 0.000 & 0.064 \\
\midrule
$P(\text{resolve})$ Rubin, $\mathcal{D}=5$ [\%]  & 16.71  & 0.0015 \\
$P(\text{resolve})$ Roman, $\mathcal{D}=5$ [\%]  & 41.63  & 6.13 \\
$P(\text{resolve})$ Roman, $\mathcal{D}=20$ [\%] & 14.02  & 0.171 \\
$P(\text{resolve})$ Roman, PSF FWHM [\%]         &  0.875 & 0.0001 \\
\bottomrule
\end{tabular}
\end{table}

\section{Satellite parallax}

\begin{figure}[h]
\centering
\includegraphics[width=\textwidth]{p6_parallax.png}
\caption{Survival function of the gain from placing Roman at L2 rather than at Earth, for the
parallax (a) and the lens mass (b). The same event is characterised twice, differing only in the
observer's position, so this is a genuinely controlled comparison. Note the logarithmic vertical
axis: the median gain is exactly $1$, and the entire effect lives in a tail of order $0.1\%$.}
\label{fig:parallax}
\end{figure}

This is a null result and worth stating plainly. The $0.01$\,au displacement of L2 gives a
parallax baseline that is negligible against the Earth's own annual motion for events lasting
months to years, which is what both of these populations produce. Fewer than $0.2\%$ of events
gain even a $10\%$ improvement in $\sigma(\piE)$. Satellite parallax is a technique for
\emph{short} events, and these are not short.

\section{Resolving the two images}

\begin{figure}[h]
\centering
\includegraphics[width=\textwidth]{p6_resolution.png}
\caption{(a) Probability that the two lensing-induced images are resolvable, for three
resolution criteria. (b) The same probability at the most permissive criterion, against lens
mass.}
\label{fig:resolution}
\end{figure}

A point lens forms two images separated by $\Delta\theta(u) = \tetE\sqrt{u^2+4}$. Following
Sajadian \& Makler (2026) we count \emph{epochs} at which both images are individually above the
single-visit depth and their separation exceeds a stated bar, and call the pair resolvable when
at least three epochs qualify. Counting epochs rather than evaluating the separation once per
event is not a refinement but a necessity: the separation is \emph{smallest} at closest approach
and grows as the source departs, while the minor image's magnification
$A_- = (u^2+2)/(2u\sqrt{u^2+4}) - 1/2$ collapses on the same motion, so the images are least
separated exactly when both are bright.

The result is dominated by the choice of criterion. With $\sigma_r = \mathcal{D}\sigma_a$ and
$\mathcal{D}=5$ -- the faint-limit value -- Roman resolves $41.6\%$ of its black-hole detections;
at $\mathcal{D}=20$ this falls to $14.0\%$, and at the diffraction-limited PSF FWHM to $0.87\%$.
Any single number here is a choice of convention presented as a result, which is why all three
are reported. Roman outperforms Rubin by roughly $2.5\times$ at a common $\mathcal{D}$ and by
three orders of magnitude at the PSF bar, reflecting its $105$\,mas F146 PSF against Rubin's
arcsecond seeing disc.

\section{Where the measurable lenses are}

\begin{figure}[h]
\centering
\includegraphics[width=\textwidth]{p7_mass_distance.png}
\caption{Weighted fraction of detected events whose lens mass is measured to better than $10\%$,
in the mass--distance plane. Measurability concentrates on \emph{nearby, massive} lenses, as
$\tetE = \sqrt{\kappa \Ml \pi_{\rm rel}}$ requires.}
\label{fig:massdist}
\end{figure}

A mass precision quoted without saying which lenses it applies to is a statement about a sample
rather than about the Galaxy. Across all detections the mass is measured to $10\%$ for $1.50\%$
of black-hole events (median lens distance $5.15$\,kpc) and $0.34\%$ of neutron-star events
(median $3.66$\,kpc), and Fig.~\ref{fig:massdist} shows those events are drawn overwhelmingly
from lenses within a few kpc.

\begin{figure}[h]
\centering
\includegraphics[width=\textwidth]{p7_sky.png}
\caption{(a) Relative weighted event yield across the scanned region, for the black-hole run.
(b) Weighted fraction of detections with $\sigma(\tE)/\tE < 0.1$ from the joint fit. The yield
concentrates toward the inner bulge; the precision is set by Rubin's visit density and is
comparatively flat.}
\label{fig:sky}
\end{figure}

\section{Caveats}

\begin{itemize}
\item \textbf{The populations are run separately and must not be pooled.} The event-rate weight
      contains $\sqrt{\Ml}$, which is valid only for the mass function actually sampled. Any
      combined black-hole plus neutron-star statement requires an assumed relative abundance,
      applied after the fact.
\item \textbf{The log-uniform black-hole mass function is an assumption, not a measurement.} It
      is taken as the truth here rather than as an importance-sampling device.
\item \textbf{No absolute yield is quoted.} The weight determines fractions, medians and ratios;
      the normalising constants cancel and are not tracked, so nothing here predicts an event
      count for either survey.
\item \textbf{The footprint sample is small.} Every Roman-versus-joint comparison rests on
      $\Neff \approx 5400$, against $\Neff \approx 1.4$--$3.0\times10^5$ for the full scan.
\item \textbf{$\sigma_{\rm joint} \le \sigma_{\rm single}$ is violated on six events} by a factor
      $1.0003$, consistent with round-off in the $9\times9$ matrix inversion on ill-conditioned
      events rather than with an error in the information sum.
\item \textbf{Rubin's astrometric error model is the weakest input.} The delivered per-visit
      floor is taken as $10$\,mas; published estimates range from $3$ to $7$\,mas, which would
      make Rubin better than assumed here and would raise its $\tetE$ column specifically.
\end{itemize}

%\section*{Reproducing this report}
%
%\begin{verbatim}
%.roman/bin/python analysis/p6_synergy_resolution.py \
%    --run bh=runs/prod_bh_20260917 --run ns=runs/prod_ns_20260917 \
%    -o figures/prod_20260918/p6
%.roman/bin/python analysis/p7_forecast_figures.py \
%    --run bh=runs/prod_bh_20260917 --run ns=runs/prod_ns_20260917 \
%    -o figures/prod_20260918/p7
%cd Report && latexmk -pdf populations_report.tex
%\end{verbatim}

\end{document}
